\documentclass[12pt]{article}
\usepackage[T1]{fontenc}
\usepackage[utf8]{inputenc}
\usepackage{lmodern}
\usepackage{textcomp}
\usepackage{enumitem}
\usepackage{geometry}
\usepackage{graphicx}
\usepackage{amsmath, amssymb}
\geometry{margin=1in}
\begin{document}
\begin{center}
Chemistry - Graduate Level Assessment 18 \
Total Marks: 40 \
Answer all questions.
\end{center}

\section*{Multiple Choice (10 marks)}
\begin{enumerate}[label=\arabic*.]
\item The symbol for antimony is
\begin{enumerate}[label=(\alph*)]
    \item Am
    \item As
    \item Ab
    \item Au
    \item Fe
\end{enumerate}
\item A saturated solution of CaF\_2 contains .00168 g of CaF\_2 per 100 g of water. Determine theK\_sp.
\begin{enumerate}[label=(\alph*)]
    \item 2.11 \ensuremath{\times} 10^-8
    \item 2.73 \ensuremath{\times} 10^-8
    \item 1.08 \ensuremath{\times} 10^-8
    \item 3.92 \ensuremath{\times} 10^-8
    \item 4.77 \ensuremath{\times} 10^-8
\end{enumerate}
\item Write theorbitals, in the order of increasing energy, for hydrogen - like atoms.
\begin{enumerate}[label=(\alph*)]
    \item 1 s \textless{} 2 p \textless{} 2 s \textless{} 3 p \textless{} 3 s \textless{} 3 d
    \item 2 s \textless{} 2 p \textless{} 1 s \textless{} 3 s \textless{} 3 p \textless{} 3 d
    \item 1 s \textless{} 2 s = 2 p \textless{} 3 s = 3 p = 3 d \textless{} 4 s = 4 p = 4 d = 4 f \textless{} 5 s = 5 p = 5 d = 5 f = 5 g
    \item 2 s = 2 p \textless{} 1 s = 3 p = 3 d
    \item 1 s \textless{} 3 s \textless{} 2 s \textless{} 5 s \textless{} 4 s
\end{enumerate}
\item Use the van der Waals parameters for chlorine to calculate approximate values of the Boyle temperature of chlorine.
\begin{enumerate}[label=(\alph*)]
    \item $1.25 \times 10^3$ $\mathrm{K}$
    \item $1.20 \times 10^3$ $\mathrm{K}$
    \item $1.41 \times 10^3$ $\mathrm{K}$
    \item $1.10 \times 10^3$ $\mathrm{K}$
    \item $1.65 \times 10^3$ $\mathrm{K}$
\end{enumerate}
\item The +1 oxidation state is more stable than the +3 oxidation state for which group 13 element?
\begin{enumerate}[label=(\alph*)]
    \item Al
    \item Tl
    \item C
    \item B
    \item Si
\end{enumerate}
\end{enumerate}

\section*{True / False (10 marks)}
\textit{Answer True or False}
\begin{enumerate}[label=\arabic*.]
\setcounter{enumi}{5}
\item True or False: The equation $x^2$ - 5 x + 6 = 0 has solutions of x = 4 and x = 5 when solved using the quadratic formula.
\item True or False: Methylamine (CH 3 NH 2) can form hydrogen bonds due to the presence of a nitrogen atom with a lone pair of electrons.
\item True or False: The 13 C spectrum of 3-ethylpentane exhibits five distinct chemical shifts due to its unique molecular structure.
\item True or False: Covalent network bonding typically results in high melting points and poor electrical conductivity in substances.
\item True or False: The logarithm of $3^2$ is equal to 3, representing the exponent directly without any calculations.
\end{enumerate}

\section*{Long Form Response (20 marks)}
\textit{Answer each question with a clear and complete explanation.}
\begin{enumerate}[label=\arabic*.]
\setcounter{enumi}{10}
\item Discuss the factors that influence the chemical shift in NMR spectroscopy and analyze why CH 3 Cl exhibits a 1 H resonance approximately 1.55 kHz from TMS in a 12.0 T magnetic field.
\item Discuss the process by which cobalt-60 is produced from cobalt-59, emphasizing the role of neutron bombardment in radiation therapy applications for cancer treatment.
\end{enumerate}

\vfill
\noindent\textit{End of Paper}
\end{document}